% !PRINTED 587
\weberpage{587}
Es ist noch zu zeigen, daß von den Zahlen (13) keine zwei
einander gleich sind. Wäre
\[
\boldsymbol{s}\left(\frac{2\varrho}{\mu}\right) = \boldsymbol{s}\left(\frac{2\varrho'}{\mu}\right),
\]
ohne daß $\varrho \equiv \varrho'$ (mod $\mu$) wäre, so müßte nach (5)
\[
\frac{2(\varrho + \varrho')}{\mu} \equiv 1 \quad (\mathrm{mod}\ 2, \omega)
\]
sein. Es müßte also $a(1+\omega)$ durch 2 teilbar sein, also auch,
wenn $\omega$ zu $\omega'$ konjugiert ist, $a(1+\omega')$ durch 2 teilbar, und
mithin
\[
a(\omega - \omega') = \sqrt{\Delta} \equiv 0 \quad (\mathrm{mod}\ 2).
\]
Es wäre also $\Delta$ durch 4 teilbar, und $\frac{1}{4}\Delta = \frac{1}{4}a^2(\omega - \omega')^2$
wäre gleichfalls noch Diskriminante, was der Annahme widerspricht,
daß $\Delta$ Stammdiskriminante sei\footnote{Es müßte $b$ gerade und $a-b+c$ durch 4 teilbar sein, und
$\frac{1}{4}\Delta = \left(\frac{b}{2}-a\right)^2 - a(a-b+c)$ wäre Diskriminante. Für eine beliebige
Diskriminante $D$ könnte man denselben Zweck erreichen, wenn man zum
Zähler der Ausdrücke (13) eine Potenz von 2 als Faktor hinzufügte.}. Hiernach sind die
Größen (13), deren Anzahl $m-1$ beträgt (Bd. II, § 165), alle
voneinander verschieden, und es ergibt sich
\weq{14}{A(x) = \overset{\varrho}{\prod}\left[x - \boldsymbol{s}\left(\frac{2\varrho}{\mu}\right)\right].}

Die Funktion $\boldsymbol{s}(u)$ genügt nach § 44, (20) der Gleichung:
\[
\boldsymbol{s}\left(u + \frac{\omega}{2}\right) = \frac{1}{\boldsymbol{s}(u)},
\]
und wegen (8) und (4):
\[
\boldsymbol{s}\left[\mu\left(u + \frac{\omega}{2}\right)\right] = \frac{(-1)^{cx}}{\boldsymbol{s}(\mu u)},
\]
und demnach ergibt sich aus (10), wenn man $u$ durch $u + \dfrac{\omega}{2}$
ersetzt:
\weq{15}{\frac{(-1)^{cx}}{\boldsymbol{s}(\mu u)} = \frac{1}{\boldsymbol{s}(u)}\,\frac{A\left(\dfrac{1}{\boldsymbol{s}}\right)}{D\left(\dfrac{1}{\boldsymbol{s}}\right)} = \frac{(-1)^{cx} D(\boldsymbol{s})}{\boldsymbol{s}(u) A(\boldsymbol{s})},}
\weq{16}{A(\boldsymbol{s}) A\left(\frac{1}{\boldsymbol{s}}\right) = (-1)^{cx} D(\boldsymbol{s}) D\left(\frac{1}{\boldsymbol{s}}\right).}
% [source sic: equations (15)--(16) behave as though the external factor \mu in (10) has been absorbed into a normalized A]
